
\section{Related Work}\label{sec:related}

\paragraph{Runtime control of speculative decoding.}
Speculative decoding accelerates autoregressive generation by verifying draft tokens in parallel while preserving the target model's output distribution \citep{leviathan2023fast,chen2023accelerating,miao2024specinfer}. Draft architectures such as Medusa and EAGLE-3 improve the efficiency of token proposals \citep{cai2024medusa,li2024eagle,li2024eagle2,li2025eagle3}. Beyond draft quality, serving performance depends on when and how speculation is used. TurboSpec \citep{liu2024turbospec} selects a speculation length, including zero, based on estimated goodput; AdaSpec \citep{huang2025adaspec} accounts for latency SLOs; Nightjar \citep{li2025nightjar} disables unprofitable speculation and offloads the drafter; and FASER \citep{chen2026faser} adjusts speculation length per request. MoESD \citep{huang2025moesd} and MagicDec \citep{sadhukhan2025magicdec} further show that the profitability of speculation depends on target architecture, batch size, and sequence length. These methods optimize speculative execution given a drafter. Our setting adds a second decision: whether spending resources to improve that drafter will increase future serving throughput enough to justify the cost.

\paragraph{Online drafter adaptation.}
Online speculative decoding addresses draft--target mismatch by updating the drafter from serving traffic. OSD \citep{liu2024osd} uses online distillation to adapt to changing query distributions, while OnlineSPEC \citep{qian2026drafts} develops online-learning algorithms with dynamic-regret guarantees. DVI \citep{bhansali2025dvi} learns from acceptance and rejection signals, and related work explores online drafter selection \citep{liu2026notabandit} and co-training during reinforcement-learning rollouts \citep{wang2026cotraining}. These approaches primarily improve how a drafter learns from available feedback. Our controller treats the learning process itself as a serving decision: higher acceptance is useful only when the resulting throughput gain justifies the resources consumed by adaptation.

\paragraph{Joint training and serving: Aurora and TIDE.}
Aurora \citep{wang2026aurora} closes the training--serving loop through asynchronous reinforcement learning from live inference traces, using a separate training server and hot-swapping updated drafters into serving. Its emphasis is on rapid adaptation, feedback efficiency, and deployment without extensive offline draft training. TIDE \citep{park2026tide} goes further in runtime coordination: it reuses target-model hidden states for training, enables speculation according to batch-dependent profitability, and selectively resumes training when acceptance drifts, retaining updates that improve acceptance. Its deployment separates inference and training across GPU resources, including heterogeneous accelerators. Both systems demonstrate that drafter training and speculative serving can operate together, and TIDE already controls when to speculate and when to adapt. Our work addresses a different resource constraint. Training shares the serving GPU, so each update directly reduces the capacity available to process requests. We therefore estimate not only whether speculation is profitable at the current load, but also whether further adaptation can improve serving enough to repay its interference cost over the expected workload duration. Per-domain learning estimates help identify cases where improvement is insufficient, while the payback calculation limits training even when the drafter can continue to improve.

\paragraph{Co-located training and inference.}
Resource sharing between training and inference has also been studied outside speculative decoding. FlexLLM \citep{oliaro2024flexllm} co-schedules inference with fine-tuning workloads under latency constraints. Our setting differs because training changes the model used by the serving process itself: allocating more GPU time to training immediately reduces serving capacity but may improve subsequent speculative decoding. The controller makes this trade-off explicitly, allocating training from available slack under light load and evaluating its expected net benefit under saturation.
